\section{Conclusion}
In this paper, we proposed MagicMotion, a trajectory-controlled image-to-video generation method that uses a ControlNet-like architecture to integrate trajectory information into the diffusion transformer. We employed a progressive training strategy, allowing MagicMotion to support three levels of trajectory control: dense masks, bounding boxes and sparse boxes. We also utilized Latent Segment Loss to enhance the ability of the model to perceive fine-grained object shapes when only provided with sparse trajectory conditions.
Additionally, we presented MagicData, a high-quality annotated dataset for trajectory-controlled video generation, created through a robust data pipeline.
Finally, we introduced MagicBench, a large-scale benchmark for evaluating trajectory-controlled video generation. MagicBench not only assessed video quality and trajectory accuracy but also took the number of controlled objects into account.
Extensive experiments on both MagicBench and DAVIS demonstrated the superiority of MagicMotion compared to previous works.
\\
\textbf{Acknowledge} This work was supported in part by the National Natural Science Foundation of China (Grant 62032006 and Grant 62472098).